% ===================================================================
%  BAGAN No. 4 — Masa Dewasa dan Masa Tua.
%
%  Traduction, faite en 2026, en indonesien d'aujourd'hui, sur l'ido
%  de texto/io. Regles de la colonne : voir 10-bagan-01.tex. Deux
%  scenes, deux numerotations : les cles portent c1 et c2.
%
%  LA PARENTE INDONESIENNE COUPE SUR UN AXE QUE LE RELEVE N'AVAIT PAS
%  ENCORE RENCONTRE. Le coreen distingue le frere de la sœur et les
%  aines des cadets ; le tamoul distingue le beau-pere du cote du mari
%  de celui du cote de l'epouse ; l'ourdou fait les deux et ajoute
%  l'oncle aine et l'oncle cadet. L'indonesien ne genre RIEN — kakak
%  est l'aine ou l'ainee, adik le cadet ou la cadette, ipar le
%  beau-frere ou la belle-sœur, mertua le beau-pere ou la belle-mere,
%  menantu le gendre ou la bru — mais il tranche l'age partout. Il
%  faut donc ajouter « laki-laki » ou « perempuan » a presque chaque
%  terme de ce tableau, et c'est la colonne la plus lourde du releve
%  sur ce point.
%
%  ET IL A UN MOT QUE PAS UNE LANGUE DU RELEVE NE POSSEDE : besan.
%  C'est la relation entre les DEUX COUPLES DE PARENTS d'un homme et
%  d'une femme maries — le pere de Ioannes est le besan du pere de
%  Martha. L'ido n'a rien et doit ecrire « la familii di la gespozi » ;
%  le francais n'a rien non plus. Or c'est exactement la scene du
%  tableau : deux familles autour d'une meme table a cause d'un
%  mariage. Une langue a un mot pour le sujet du tableau, et ni la
%  langue construite ni la langue source ne l'ont.
%
%  beranda (8) EST L'AUTRE BOUT DU VOYAGE DE برآمدہ. La colonne
%  ourdoue avait montre que l'anglais « verandah » vient de l'ourdou,
%  et non l'inverse. L'indonesien dit beranda, qui lui vient du
%  portugais varanda — l'autre extremite du meme mot, arrive par la
%  mer et non par la terre. Deux colonnes du releve nomment cette
%  galerie par un mot qui a fait le tour du monde en sens contraire.
%
%  UN TROU FRANC, ET ON LE DIT PLUTOT QUE DE LE MASQUER. La perdrix
%  (2) n'a pas de nom indonesien : l'oiseau ne vit pas dans
%  l'archipel, et « puyuh » est la caille, « pegar » le faisan. La
%  colonne ecrit « ayam hutan », qui decrit un gallinace sauvage —
%  c'est une periphrase, comme le « grand canard » tamoul pour l'oie
%  et les « batons a battre » de cinq colonnes pour le fleau. On ne
%  met pas un autre oiseau a la place.
%
%  ET LA MALADIE RAMENE LE NEERLANDAIS : apotek la pharmacie (16),
%  dokter le medecin (25), resep l'ordonnance (27), rematik. Le
%  tableau 2 avait le corps entierement malais ; ici, dès qu'il s'agit
%  de SOIGNER ce corps, la couche coloniale reparait entiere. Le
%  partage tient toujours : ce qu'on a est malais, ce qu'on achete et
%  ce qu'on apprend a l'ecole est neerlandais.
% ===================================================================

%%K t04-apar-1 apar
\VUsaut{4.0mm}
\VUtitre{15.0pt}{60}{BAGAN No. 4}
\VUsaut{1.6mm}
\VUtitre{11.4pt}{0}{Masa dewasa dan masa tua.}
\VUsaut{3.2mm}

%%K t04-tit-1 sub
\VUcentre{11.4pt}{0}{\textit{Adegan Pertama.}}
\VUsaut{1.6mm}
\VUtitre{11.4pt}{0}{Pesta pernikahan.}
\VUsaut{2.6mm}

%%K t04-tit-2 sub
\VUpk{10.2pt}{40}{Musim gugur.}
\VUsaut{3.0mm}

%%K t04-c1-01-1 p
1. --- Kaum muda telah tumbuh besar. Salah seorang dari mereka,
Ioannes, menikah. Kita menghadiri \VUgras{pesta pernikahan} yang
mengumpulkan, di sekeliling meja yang sama, keluarga kedua mempelai.

%%K t04-c1-02-1 p
2. --- Kita berada pada \VUgras{musim gugur,} di bulan
\VUgras{September.} Angin dingin bulan \VUgras{Oktober} dan
\VUgras{November} belum melucuti pedesaan dari dedaunannya yang hijau.
Udara sore hangat dan harum ; karena itu makan malam berlangsung di
bawah \VUgras{beranda} \textsuperscript{(8)} yang menghadap kebun.

%%K t04-c1-03-1 p
3. --- Jauh di sana, di atas \VUgras{bukit} \textsuperscript{(3)},
terletak rumah orang tua Ioannes, sebuah \VUgras{kastil}
\textsuperscript{(1)} yang anggun tersembunyi di antara pepohonan ;
kebun-kebun anggur yang bagus, yang menghasilkan anggur yang sangat
baik, menutupi lereng bukit itu. Petani anggur menggarap dan
merawatnya.

%%K t04-c1-04-1 p
4. --- Di atas kebun anggur itu lewat sekawanan \VUgras{ayam hutan}
\textsuperscript{(2)} yang ketakutan oleh mendekatnya \VUgras{penjaga
hutan} \textsuperscript{(5)}. Orang itu, dengan \VUgras{senapan} di
bawah lengan dan \VUgras{tas buruan} tersandang, memeriksa
\VUgras{semak belukar} \textsuperscript{(4)} bersama \VUgras{anjingnya}
\textsuperscript{(6)}. Ia menjaga tanah itu dan mencegah para pemburu
gelap membinasakan hewan buruan.

%%K t04-c1-05-1 p
5. --- Di luar \VUgras{beranda} memanjat sebuah \VUgras{para-para
anggur} yang darinya menggantung \VUgras{tandan-tandan anggur}
\textsuperscript{(7)} yang lebat.

%%K t04-c1-06-1 p
6. --- Bunga-bunga musim gugur yang indah, yang menghiasi
\VUgras{meja} \textsuperscript{(16)}, adalah \VUgras{bunga seruni}
\textsuperscript{(9)} ; \VUgras{ayah} \textsuperscript{(10)} mempelai
perempuan menyematkan salah satunya di lubang kancing jasnya.

%%K t04-c1-07-1 p
7. --- Santapan hampir selesai. \VUgras{Pelayan} \textsuperscript{(11)}
mulai membawa hidangan penutup dan sebentar lagi akan mengangkat
bagian-bagian alat makan yang tidak diperlukan lagi :
\VUgras{sendok} \textsuperscript{(14)}, \VUgras{garpu}
\textsuperscript{(12)}, \VUgras{pisau} \textsuperscript{(13)} dan
\VUgras{pinggan besar} \textsuperscript{(15)}.

%%K t04-tit-3 sub
\VUpk{10.2pt}{40}{Sanak saudara.}
\VUsaut{3.0mm}

%%K t04-c1-08-1 p
8. --- Semua orang tampak gembira, dan senyum \VUgras{ibu}
\textsuperscript{(17)} mempelai laki-laki ketika memandang anak-anaknya
membuktikan kebahagiaannya.

%%K t04-c1-09-1 p
9. --- Pernikahan itu menyatukan dua keluarga kaya di daerah itu.
Ioannes, \VUgras{mempelai laki-laki} \textsuperscript{(18)}, adalah
\VUgras{anak} seorang tuan tanah besar, dan ia menjadi
\VUgras{menantu laki-laki} seorang pengusaha yang cakap.

%%K t04-c1-10-1 p
10. --- Kedua \VUgras{mempelai} masih muda, namun mereka sudah lama
\VUgras{bertunangan.} \VUgras{Perempuan muda} \textsuperscript{(19)}
itu, Martha, tampak memesona dalam \VUgras{gaun putihnya} yang dihiasi
bunga jeruk.

%%K t04-c1-11-1 p
11. --- \VUgras{Ayah mertuanya} \textsuperscript{(20)} sangat berbahagia
mempunyai \VUgras{menantu perempuan} yang begitu ramah, dan \VUgras{ibu
mertua} \textsuperscript{(21)} Ioannes tersenyum melihat \VUgras{anak
perempuannya} begitu cantik.

%%K t04-c1-12-1 p
12. --- Di ujung meja duduk \VUgras{para tamu} \textsuperscript{(22)}
yang lain : \VUgras{kerabat, teman, sepupu laki-laki, sepupu
perempuan.} Seorang \VUgras{kepala pelayan} \textsuperscript{(23)},
dengan serbet di bawah lengan, melayani mereka dengan rajin.

%%K t04-tit-4 sub
\VUpk{10.2pt}{40}{Para sahabat.}
\VUsaut{3.0mm}

%%K t04-c1-13-1 p
13. --- Di sebelah kanan meja, inilah seorang \VUgras{nona}
\textsuperscript{(24)}, \VUgras{teman} mempelai perempuan, lalu
\VUgras{saudara laki-lakinya}, yang menjadi \VUgras{pengiring
pengantin laki-laki} \textsuperscript{(25)}. Ia berbincang dengan
\VUgras{pengiring pengantin perempuannya} \textsuperscript{(26)},
saudara perempuan mempelai laki-laki, yang dengan pernikahan itu
menjadi \VUgras{ipar perempuan} Martha. Ia seorang gadis yang
memesona. Ia memakai \VUgras{perhiasan} yang indah, \VUgras{kalung
mutiara} dan \VUgras{gelang} emas.

%%K t04-c1-14-1 p
14. --- Di dekat mereka, tuan yang berdiri dan mengangkat gelasnya
adalah seorang \VUgras{teman} \textsuperscript{(27)} keluarga itu. Ia
salah seorang \VUgras{saksi mempelai laki-laki.} Ia
\VUgras{bersulang,} minum untuk \VUgras{kesehatan} kedua mempelai muda
dan mengharapkan mereka berbahagia lama. Ia berpakaian setelan
upacara, jas berekor dengan \VUgras{sepatu pernis}
\textsuperscript{(28)}. Ia pengiring \VUgras{nenek}
\textsuperscript{(29)}, yang sudah \VUgras{menjanda.}

%%K t04-c1-15-1 p
15. --- Pemuda yang memakai \VUgras{jas berekor}
\textsuperscript{(31)}, dengan kumis hitam tipis, adalah \VUgras{ipar
laki-laki} \textsuperscript{(30)} Ioannes. Ia seorang insinyur
terkemuka. Ia duduk di sebelah seorang \VUgras{nyonya muda}
\textsuperscript{(33)} yang sangat anggun, \VUgras{kakak perempuan}
Martha dan ibu anak perempuan kecil yang menghampiri, sambil
menggandeng sepupu laki-lakinya, untuk mempersembahkan bunga dan
\VUgras{mengucapkan} kepadanya \VUgras{selamat malam.} Nyonya itu
memakai \VUgras{anting-anting} yang sangat indah dan \VUgras{tutup
kepala} yang anggun.

%%K t04-c1-16-1 p
16. --- Sepupu laki-laki kecil itu, yang tampak begitu aneh dengan
\VUgras{blusnya} \textsuperscript{(36)} dan \VUgras{celana pendeknya}
\textsuperscript{(37)} yang menggembung, sangat patut dikasihani. Ia
\VUgras{anak yatim piatu} \textsuperscript{(35)}. Ketika masih sangat
kecil ia kehilangan ayah dan ibunya. Pamannya menjadi \VUgras{walinya}
\textsuperscript{(38)}, dan tetap membujang untuk mendidik anak yang
malang itu. Ia lelaki yang baik hati. Ia agak botak dan sangat rabun ;
ia memakai \VUgras{kacamata jepit.}

%%K t04-tit-5 sub
\VUsaut{2.4mm}
\VUpk{10.2pt}{40}{Hidangan penutup.}
\VUsaut{3.0mm}

%%K t04-c1-17-1 p
17. --- Di belakang para tamu, di atas meja yang ditutup
\VUgras{taplak meja,} \VUgras{hidangan penutup} \textsuperscript{(39)}
telah disiapkan. Di dalam mangkuk besar, inilah buah-buahan :
\VUgras{buah persik} \textsuperscript{(40)}, \VUgras{buah pir}
\textsuperscript{(41)}, \VUgras{apel} \textsuperscript{(42)},
\VUgras{aprikot} \textsuperscript{(43)} dan \VUgras{anggur}
\textsuperscript{(44)}. Di dekatnya, \VUgras{nanas}
\textsuperscript{(45)}, \VUgras{kue} \textsuperscript{(46)} yang indah,
\VUgras{botol-botol minuman keras} \textsuperscript{(48)} dan
\VUgras{sampanye} \textsuperscript{(49)}, serta tumpukan
\VUgras{piring} \textsuperscript{(47)} yang bersih.

%%K t04-c1-18-1 p
18. --- \VUgras{Juru gudang anggur} \textsuperscript{(50)} membuka
sebuah \VUgras{botol} \textsuperscript{(52)} anggur tua dengan
\VUgras{pembuka botol} \textsuperscript{(51)}. \VUgras{Sumbat gabusnya}
\textsuperscript{(53)} tampak sangat sukar ditarik. Juru gudang itu
membawa \VUgras{serbet} \textsuperscript{(54)} di bawah lengannya.

%%K t04-c1-19-1 p
19. --- Sesudah makan malam, para tuan akan pergi ke ruang merokok,
dan para nyonya akan bersiap-siap untuk pesta dansa.

%%K t04-tit-6 sub
\VUsaut{5.0mm}
\VUcentre{11.4pt}{0}{\textit{Adegan Kedua.}}
\VUsaut{1.6mm}
\VUpk{11.4pt}{40}{Ulang tahun kakek.}
\VUsaut{2.6mm}

%%K t04-tit-7 sub
\VUpk{10.2pt}{40}{Kunjungan.}
\VUsaut{3.0mm}

%%K t04-c2-01-1 p
1. --- Sepuluh tahun berlalu ; orang tua Ioannes menjadi tua dan
sangat menyayangi ketiga cucu mereka, dua \VUgras{cucu laki-laki}
\textsuperscript{(2)} dan seorang \VUgras{cucu perempuan}
\textsuperscript{(3)}. Karena hari ini \VUgras{ulang tahun kakek,}
anak-anak mereka datang mengucapkan selamat.

%%K t04-c2-02-1 p
2. --- Si kecil Petrus masih sangat kecil ; umurnya baru tiga tahun.
Ia melihat \VUgras{kucing} \textsuperscript{(1)} mendekat. Ia hendak
membelai \VUgras{bulunya} yang seperti sutra dan \VUgras{ekornya}
yang panjang ; ia tidak berpikir bahwa di bawah \VUgras{kaki} yang
anggun itu tersembunyi cakar-cakar runcing yang jahat, yang barangkali
akan mencakarnya.

%%K t04-c2-03-1 p
3. --- Kakak perempuannya, Gabriela, yang menggandeng tangannya, jauh
lebih bersungguh-sungguh, dan Marcellus dengan \VUgras{mantelnya}
\textsuperscript{(4)} yang besar dan \VUgras{ikat pinggang}
\textsuperscript{(5)} kulit tampak agak kelelakian, meskipun celana
pendeknya memperlihatkan \VUgras{kaus kakinya} \textsuperscript{(6)}.

%%K t04-c2-04-1 p
4. --- Ayah mereka memakai \VUgras{mantel luar} \textsuperscript{(7)}
musim dinginnya yang tebal dengan \VUgras{kerah bulu}
\textsuperscript{(8)}, dan di dalam \VUgras{sakunya}
\textsuperscript{(9)} ada selendang. Istrinya memakai topi laken yang
dihiasi \VUgras{bulu burung} \textsuperscript{(10)}, \VUgras{jaketnya}
\textsuperscript{(11)}, \VUgras{selendang bulunya}
\textsuperscript{(12)} dan \VUgras{sarung tangan bulunya}
\textsuperscript{(13)}.

%%K t04-c2-05-1 p
5. --- Udara sangat dingin, sebab musim dingin sedang berkuasa.
\VUgras{Salju} \textsuperscript{(19)} turun sepanjang hari dan atap
rumah-rumah seluruhnya putih. Bersamaan dengan \VUgras{malam,}
dinginnya menjadi lebih menusuk lagi. Di langit, \VUgras{bulan}
\textsuperscript{(17)} dan \VUgras{bintang-bintang}
\textsuperscript{(18)} bersinar dan menembus \VUgras{kegelapan.}

%%K t04-tit-8 sub
\VUsaut{4.2mm}
\VUpk{10.2pt}{40}{Penyakit.}
\VUsaut{3.0mm}

%%K t04-c2-06-1 p
6. --- \VUgras{Orang buta} \textsuperscript{(20)} yang malang, yang
melintasi alun-alun di depan \VUgras{apotek} \textsuperscript{(16)}
dituntun oleh \VUgras{anjing pudelnya} \textsuperscript{(21)}, sangat
tidak beruntung. Iustinia, \VUgras{pembantu perempuan}
\textsuperscript{(14)}, membawa dari dapur sebuah \VUgras{mangkuk}
\textsuperscript{(15)} berisi \VUgras{air rebusan daun} yang sangat
panas dan diberi banyak gula.

%%K t04-c2-07-1 p
7. --- \VUgras{Nenek} \textsuperscript{(23)}, yang hendak mencari
\VUgras{kacamata bertangkainya} \textsuperscript{(26)} di atas meja
untuk membaca \VUgras{resep} \textsuperscript{(27)} yang baru saja
ditulis \VUgras{dokter} \textsuperscript{(25)}, bangkit dari kursi
besarnya sambil bertumpu pada \VUgras{tongkatnya}
\textsuperscript{(24)}.

%%K t04-c2-08-1 p
8. --- Sering ia menderita \VUgras{penyakit-penyakit kecil, pusing,
pilek, sakit gigi} ; kakinya lemah dan matanya tidak begitu baik lagi.
Meskipun demikian, dengan senang hati ia mengolok-olok \VUgras{dokter}
tuanya yang selalu hendak meresepkan \VUgras{obat cair}
\textsuperscript{(28)} kepadanya. Hari ini ia sangat gembira karena
seluruh keluarga mudanya berada di sekelilingnya.

%%K t04-c2-09-1 p
9. --- Terasa nyaman di dalam kamar yang tertutup rapat itu. Di
\VUgras{perapian} \textsuperscript{(29)} marmer menyala \VUgras{api
yang indah} \textsuperscript{(30)}, dan orang merasa berbahagia dalam
\VUgras{cahaya} lembut \VUgras{lampu} \textsuperscript{(31)} yang baru
saja dinyalakan.

%%K t04-tit-9 sub
\VUsaut{4.2mm}
\VUpk{10.2pt}{40}{Kakek.}
\VUsaut{3.0mm}

%%K t04-c2-10-1 p
10. --- Bahkan \VUgras{kakek} \textsuperscript{(33)} pun lupa bahwa ia
sakit, bahwa \VUgras{rematiknya} mengikatnya pada \VUgras{kursi
besarnya} \textsuperscript{(34)}, dan bahwa kemarin pun ia masih
\VUgras{demam dan pingsan.} Ia tidak lagi ingat bahwa pada musim
dingin yang lalu ia menderita \VUgras{radang paru-paru} dan bahwa
\VUgras{batuk} yang serak dan menyiksa membuatnya sangat menderita.

%%K t04-c2-10-2 p
Namun ia sembuh dengan sangat susah payah. Sekarang keadaannya agak
lebih baik. Tetapi kakinya, yang ia sandarkan pada sebuah
\VUgras{bantal} \textsuperscript{(38)} di atas \VUgras{bangku kecil}
\textsuperscript{(39)}, masih juga terasa sakit.

%%K t04-c2-11-1 p
11. --- Bila ia terbungkus rapi dalam \VUgras{jubah kamarnya}
\textsuperscript{(35)} yang hangat dengan \VUgras{kancing-kancing}
\textsuperscript{(36)} tebal dari mutiara, dan bila di dekatnya ada
\VUgras{anak yatim perempuan} \textsuperscript{(40)} kecil yang
membacakan koran untuknya — berpakaian \VUgras{tudung kepala} putih,
\VUgras{gaun} \textsuperscript{(42)} biru dan \VUgras{mantel pendek}
\textsuperscript{(41)} —, ia tidak mengeluh.

%%K t04-c2-12-1 p
12. --- Setiap tahun, ketika \VUgras{kalender} \textsuperscript{(43)}
kembali menunjukkan \VUgras{tanggal} satu \VUgras{Desember,} ia
menantikan \VUgras{cucu-cucunya} dengan tidak sabar, sebab ia tahu
bahwa mereka tidak akan melupakan \VUgras{tanggal} yang penting itu.

%%K t04-c2-13-1 p
13. --- Dengan terharu ia mengenang masa yang sudah sangat lampau,
ketika ia sendiri datang mengucapkan selamat kepada \VUgras{marsekal}
kekaisaran yang termasyhur, yang \VUgras{potretnya}
\textsuperscript{(44)} tergantung di dinding dan yang menjadi
\VUgras{buyut} anak-anaknya.

\VUsaut{6.0mm}
\VUfilet{22mm}
